\chapter{Conclusion}
\label{ch:conclusion}

Legacy software development environments suffer from configuration complexity, inconsistent build outcomes, and substantial onboarding overhead.
This thesis investigated whether containerization could modernize development workflows for native software projects without fundamental codebase changes.
Applied to an industrial SCADA system, this work produced empirical evidence that containerized environments substantially improve developer productivity and build performance compared to traditional approaches.

The prototype demonstrated that containers integrate through three principles: automated provisioning using infrastructure-as-code, seamless IDE integration via DevContainers, and standardized cross-platform build configurations.
Toolchain switching time reduced by 95\%---from 12.9 minutes to 37.5 seconds---validating direct productivity improvements for developer onboarding and multi-version development.

Performance evaluation across six environments revealed nuanced characteristics.
Linux-based containers outperformed Windows-native approaches by 35--72\% for filesystem-intensive operations, regardless of native or WSL2 execution.
For CPU-bound compilation, all configurations performed within 15\% of each other, with native Linux containers imposing minimal overhead.
Compiler caching amplified differences: Linux containers achieved 24--27$\times$ speedup versus 4--5$\times$ on Windows.
Combined optimization techniques (caching and parallel testing) reduced total workflow time by up to 89\% compared to unoptimized Windows workflows.

Resource utilization distinguished virtualization from containerization overhead.
CPU utilization reflected operation types---compilation saturated cores regardless of environment.
Hypervisor-based solutions require 5--10 GB additional host memory, making native containerization advantageous for constrained systems.

\begin{sloppypar}
This work provides quantitative benchmarks, architectural patterns, and implementation guidance for modernizing development infrastructure.
The infrastructure-as-code approach addresses the ``works on my machine'' problem while maintaining compatibility with existing systems.
It demonstrates that containerization principles from microservice deployment transfer effectively to monolithic legacy systems.
\end{sloppypar}

Future research should extend validation across diverse codebases and programming languages, complemented by qualitative studies of developer experience.
The low overhead of native containers suggests opportunities for investigating remote development infrastructure.

Containerization represents not merely a deployment technology but a fundamental improvement to legacy software development workflows.
It delivers reproducibility, productivity, and performance benefits that make it a viable modernization strategy for organizations preserving established software investments while maintaining competitiveness.
